\section{Auditoría de fuentes: convertir un tema sensible en una clase de sistemas verificable}

Esta sesión la imparte Julie Cordua, CEO de Thorn, y trata de cómo la tecnología de seguridad infantil en línea ayuda a las plataformas a identificar riesgos, reducir la difusión repetida, organizar la revisión humana y entregar las pistas que cumplen los requisitos al sistema legal de reporte e investigación. Al consultarlo el 11 de agosto de 2026, el video público pedía iniciar sesión; el repositorio conserva la portada oficial y 833 subtítulos con marcas de tiempo, por lo que este texto toma los subtítulos de la clase como hilo narrativo y contrasta los mecanismos permanentes con los materiales oficiales de Thorn, NCMEC y NIST, sin presentar capacidades de producto posteriores como hechos de la clase de 2025.

\begin{warningbox}{Tratamiento del contenido sensible y límites de la evidencia}
Este texto no relata, describe ni muestra contenido ilegal, ni ofrece métodos para generarlo, difundirlo o evadir su detección. Los casos individuales mencionados en clase se abstraen únicamente como eventos del sistema, por ejemplo: «la plataforma detectó un archivo de alto riesgo y completó el reporte». El número de reportes, el número de archivos, las proporciones de los estudios y la escala de los productos dependen del año, de la definición y del método de deduplicación; salvo que se indique de forma explícita, las cifras de años distintos no deben restarse ni compararse directamente.
\end{warningbox}

\begin{importantbox}{Los tres ciclos cerrados de esta sesión}
\begin{enumerate}
  \item \textbf{Response loop}: carga o mensaje $\rightarrow$ detección $\rightarrow$ prioridad $\rightarrow$ revisión humana $\rightarrow$ medidas de la plataforma y reporte legal;
  \item \textbf{Learning loop}: resultados verificados y autorizados $\rightarrow$ trusted data/hash/model update $\rightarrow$ cobertura de riesgos futuros;
  \item \textbf{Governance loop}: nuevas tecnologías y nuevos abusos $\rightarrow$ threat research $\rightarrow$ Safety by Design $\rightarrow$ pruebas, auditoría y transparencia.
\end{enumerate}
\end{importantbox}

\begin{knowledgebox}{Términos clave: no mezclar cuatro tipos de objeto bajo «contenido malo»}
\textbf{Content} es el archivo o mensaje que procesa la plataforma; \textbf{signal} es una hash match, un score del modelo o un patrón de comportamiento; \textbf{case} es un evento revisable organizado por cuenta, tiempo y contexto; \textbf{outcome} es el resultado en forma de medidas de la plataforma, reporte legal, investigación o protección de la víctima. Toda métrica de ingeniería debe indicar qué capa mide; de lo contrario, «aumentó el volumen de detección» puede significar una mejor cobertura, pero también solo archivos repetidos, un cambio de umbral o un cambio en el criterio de reporte.
\end{knowledgebox}

\subsection{Por qué esto no es «entrenar un clasificador»}

Esta sección define primero el objeto de ingeniería. \term{CSAM} es la sigla de child sexual abuse material y designa el material de abuso sexual infantil tal como lo definen la ley y las políticas de las plataformas; el término describe a la víctima y la evidencia, y no debe confundirse con «contenido para adultos» común. Un modelo solo puede producir una match o un score; el resultado real de protección depende además de la policy de la plataforma, la cola de moderación, la preservación de la evidencia, los organismos legales de reporte, la investigación policial y la identificación de las víctimas. Si falla cualquier eslabón, la precisión del eslabón anterior no se convierte automáticamente en un resultado de seguridad.

\lecturefigure{01-response-pipeline.png}{La tecnología de seguridad infantil es una cadena de respuesta de extremo a extremo formada por detect, review, action/report y victim identification.}{Subtítulos locales 00:00--05:50; rediseño conceptual no sensible basado en el flujo presentado en clase.}

\begin{knowledgebox}{Lectura de la figura: distinguir primero entre signal, decision y outcome}
Lo que producen el classifier o el hash service de la izquierda es una \textbf{señal}; la moderación de la plataforma y la policy dan lugar a una \textbf{decisión de producto}; los organismos de reporte y los investigadores buscan un \textbf{resultado de investigación} dentro del procedimiento legal. El score de un modelo no puede probar por sí solo una identidad, una intención subjetiva ni una responsabilidad legal. El diseño del sistema debe registrar quién ejecuta cada transición de estado, con base en qué policy, qué audit evidence deja y cómo se apelan o corrigen los errores.
\end{knowledgebox}

\teachervoice{Cordua no abre la sesión con casos impactantes para captar la atención, sino con una pregunta: cuando la carga de reportes de las plataformas no deja de crecer, ¿puede la tecnología conectar de verdad la detección, las medidas adoptadas y el tiempo limitado de los investigadores? Esta pregunta lleva la clase de la «identificación de contenido malo» a las end-to-end operations.}

\subsection{Resumen de la sección}

El objetivo de esta sesión es explicar el sistema de respuesta completo, no hacer promoción de producto para un modelo en particular. Todos los mecanismos que siguen deben responder tres preguntas: qué tipo de riesgo reducen, a quién le transfieren el trabajo y dónde están los límites de su evidencia.
